% Figure: a stepped circular shaft carrying a single torque T, fixed at one
% end. Two segments of different polar second moment store torsional strain
% energy in proportion that the worked example partitions. Drawn as a side
% elevation of the two cylinders with the applied torque arrow.
\begin{figure}[htbp]
\centering
\begin{tikzpicture}[
  shaft/.style={draw=engblack, very thick},
  fillA/.style={fill=engblue!14},
  fillB/.style={fill=engamber!18},
  torque/.style={draw=engred, -{Stealth}, very thick},
  dimline/.style={draw=enggray, {Stealth}-{Stealth}, thin},
  label/.style={font=\small},
]
% fixed wall at the left
\draw[shaft] (0,-1.1) -- (0,1.1);
\foreach \y in {-0.9,-0.5,-0.1,0.3,0.7} {
  \draw[draw=enggray, thin] (0,\y) -- (-0.28,\y+0.2);
}
\node[font=\scriptsize, anchor=south] at (-0.3,1.0) {fixed};

% segment 1 (fat, large J_1), length L_1 from x=0 to x=4
\fill[engblue!14] (0,-0.8) rectangle (4,0.8);
\draw[shaft] (0,-0.8) -- (4,-0.8);
\draw[shaft] (0,0.8) -- (4,0.8);
\draw[shaft] (0,-0.8) -- (0,0.8);
\node[engblue, label, anchor=south] at (2,0.85) {segment 1: $J_1,\ L_1$};

% segment 2 (thin, small J_2), length L_2 from x=4 to x=8
\fill[engamber!18] (4,-0.45) rectangle (8,0.45);
\draw[shaft] (4,-0.45) -- (8,-0.45);
\draw[shaft] (4,0.45) -- (8,0.45);
\draw[shaft] (4,-0.8) -- (4,-0.45);
\draw[shaft] (4,0.45) -- (4,0.8);
\draw[shaft] (8,-0.45) -- (8,0.45);
\node[engamber, label, anchor=south] at (6,0.5) {segment 2: $J_2,\ L_2$};

% the applied torque T at the free end, drawn as a curved arrow
\draw[torque] (8.55,0) arc[start angle=-80, end angle=200, radius=0.55];
\node[engred, label, anchor=west] at (9.2,0) {$T$};

% the same internal torque T runs through both segments (statics)
\node[engblack, font=\scriptsize, anchor=north] at (2,-0.95) {internal torque $T$};
\node[engblack, font=\scriptsize, anchor=north] at (6,-0.6) {internal torque $T$};

% dimension lines
\draw[dimline] (0,-1.5) -- (4,-1.5);
\node[enggray, font=\scriptsize, anchor=north] at (2,-1.52) {$L_1$};
\draw[dimline] (4,-1.5) -- (8,-1.5);
\node[enggray, font=\scriptsize, anchor=north] at (6,-1.52) {$L_2$};

\end{tikzpicture}
\caption[Stepped torsion shaft for the strain-energy partition]{A circular
shaft fixed at the left and twisted by a single torque $T$ at the free
right end, stepped from a fat segment of polar second moment $J_1$ and
length $L_1$ to a thin segment $J_2$, $L_2$. Static equilibrium carries
the same internal torque $T$ through both segments, so each stores
torsional strain energy $T^2 L_i/(2 G J_i)$, and the two energies stand in
the ratio $(L_1/J_1):(L_2/J_2)$. The thin segment, with the smaller $J$,
stores the larger share of the energy and carries the larger share of the
twist, which is why a single under-sized step dominates the compliance of
an otherwise stiff shaft. The worked example partitions the energy and
recovers the total angle of twist by Castigliano's theorem.}
\label{fig:vol03ch05:stepped-shaft}
\end{figure}
